% ============================================================
%  يُترجَم بـ XeLaTeX  (في Overleaf: Menu → Compiler → XeLaTeX)
%  الخط: Amiri (متوفر في Overleaf، أو حمّله من github.com/aliftype/amiri)
% ============================================================
\documentclass[12pt,a4paper]{article}
\usepackage[a4paper,margin=1.7cm,top=2.3cm,bottom=2.9cm,headheight=15pt,footskip=1.5cm]{geometry}
\usepackage{amsmath,amssymb}
\usepackage{icomma}
\usepackage[english,bidi=default,layout=lists]{babel}
\babelprovide[import,main]{arabic}
\babelfont{rm}[Scale=1.08]{Amiri}
\babelfont{sf}[Scale=1.08]{Amiri}
\usepackage[most]{tcolorbox}
\usepackage{tikz}
\usetikzlibrary{arrows.meta,calc,shadings}
\usepackage{array}
\usepackage{colortbl}
\usepackage{fancyhdr}
\usepackage{titlesec}
\usepackage{enumitem}
\usepackage{collcell}
\usepackage{needspace}
\usepackage{etoolbox}
\newcolumntype{C}{>{\collectcell\ar}c<{\endcollectcell}}
\pretocmd{\section}{\Needspace{15\baselineskip}}{}{}
\pretocmd{\subsection}{\Needspace{11\baselineskip}}{}{}

\definecolor{cDef}{HTML}{1F4E9C}
\definecolor{cThm}{HTML}{B3261E}
\definecolor{cRem}{HTML}{2E7D32}
\definecolor{cEx}{HTML}{C76B00}
\definecolor{cMeth}{HTML}{6A1B9A}
\definecolor{cSol}{HTML}{00695C}
\definecolor{cSec}{HTML}{1F3A68}

\titleformat{\section}{\Large\bfseries\color{cSec}}{\thesection.}{0.6em}{}[{\color{cSec!60}\titlerule[1pt]}]
\titleformat{\subsection}{\large\bfseries\color{cThm!85!black}}{\thesubsection.}{0.5em}{}
\titlespacing*{\section}{0pt}{14pt}{8pt}
\titlespacing*{\subsection}{0pt}{10pt}{5pt}
\setlength{\parindent}{0pt}
\setlength{\parskip}{4pt}
\setlist{itemsep=2pt,topsep=3pt}

\newtcolorbox{arbox}[3][]{enhanced,
  colback=#2!5,colframe=#2!80!black,arc=2mm,boxrule=0.8pt,
  left=6pt,right=6pt,top=8pt,bottom=5pt,
  fonttitle=\bfseries,coltitle=white,
  title={\foreignlanguage{arabic}{#3}},
  attach boxed title to top right={xshift=-5mm,yshift=-3mm},
  boxed title style={colback=#2!80!black,arc=1.5mm,boxrule=0pt},
  before upper={\selectlanguage{arabic}},
  before skip=8pt,after skip=8pt,breakable,#1}
\newenvironment{BOX}[2]{\begin{otherlanguage}{english}\begin{arbox}{#1}{#2}}{\end{arbox}\end{otherlanguage}}
\newenvironment{tarif}[1][تعريف]{\begin{BOX}{cDef}{#1}}{\end{BOX}}
\newenvironment{khassia}[1][خاصية]{\begin{BOX}{cThm}{#1}}{\end{BOX}}
\newenvironment{molahada}[1][ملاحظة]{\begin{BOX}{cRem}{#1}}{\end{BOX}}
\newenvironment{mithal}[1][مثال]{\begin{BOX}{cEx}{#1}}{\end{BOX}}
\newenvironment{tariqa}[1][طريقة]{\begin{BOX}{cMeth}{#1}}{\end{BOX}}

\newcounter{exo}
\newenvironment{tamrin}[1]{\stepcounter{exo}\begin{BOX}{cSec}{التمرين \theexo: #1}}{\end{BOX}}
\newenvironment{hal}{\Needspace{12\baselineskip}\begin{otherlanguage}{english}\begin{arbox}[breakable]{cSol}{حل التمرين \theexo}}{\end{arbox}\end{otherlanguage}}

\newenvironment{LTR}{\begin{otherlanguage}{english}}{\end{otherlanguage}}
\newcommand{\ar}[1]{\foreignlanguage{arabic}{#1}}

\renewcommand{\dots}{\mathinner{\ldotp\ldotp\ldotp}}
\renewcommand{\ldots}{\mathinner{\ldotp\ldotp\ldotp}}
\renewcommand{\cdots}{\mathinner{\cdotp\cdotp\cdotp}}
\newcommand{\N}{\mathbb{N}}
\newcommand{\Z}{\mathbb{Z}}
\newcommand{\D}{\mathbb{D}}
\newcommand{\Q}{\mathbb{Q}}
\newcommand{\R}{\mathbb{R}}
\newcommand{\pgcd}{\operatorname{PGCD}}
\newcommand{\ppcm}{\operatorname{PPCM}}

% \NLseg{a}{b}{[}{]} ; \NLright{a}{[} ; \NLleft{b}{]}  (interval on a number line)
\newcommand{\NLseg}[4]{\begin{LTR}\begin{tikzpicture}[baseline=-0.6ex]
 \draw[-{Stealth[length=2mm]},gray!70](0,0)--(3.4,0);
 \draw[line width=2pt,cThm!80!black](0.7,0)--(2.7,0);
 \node[cThm!80!black,font=\LARGE] at (0.7,0){$#3$};
 \node[cThm!80!black,font=\LARGE] at (2.7,0){$#4$};
 \node[below=3pt,font=\footnotesize] at (0.7,0){$#1$};
 \node[below=3pt,font=\footnotesize] at (2.7,0){$#2$};
\end{tikzpicture}\end{LTR}}
\newcommand{\NLright}[2]{\begin{LTR}\begin{tikzpicture}[baseline=-0.6ex]
 \draw[-{Stealth[length=2mm]},gray!70](0,0)--(3.4,0);
 \draw[line width=2pt,cThm!80!black](1.3,0)--(3.3,0);
 \node[cThm!80!black,font=\LARGE] at (1.3,0){$#2$};
 \node[below=3pt,font=\footnotesize] at (1.3,0){$#1$};
\end{tikzpicture}\end{LTR}}
\newcommand{\NLleft}[2]{\begin{LTR}\begin{tikzpicture}[baseline=-0.6ex]
 \draw[-{Stealth[length=2mm]},gray!70](0,0)--(3.4,0);
 \draw[line width=2pt,cThm!80!black](0.05,0)--(2.1,0);
 \node[cThm!80!black,font=\LARGE] at (2.1,0){$#2$};
 \node[below=3pt,font=\footnotesize] at (2.1,0){$#1$};
\end{tikzpicture}\end{LTR}}

% ---------- social-media icons (pure TikZ) ----------
\definecolor{cFB}{HTML}{1877F2}
\definecolor{cTG}{HTML}{229ED9}
\definecolor{cIG1}{HTML}{FEDA75}
\definecolor{cIG2}{HTML}{D62976}
\definecolor{cIG3}{HTML}{4F5BD5}
\newcommand{\fbicon}{\begin{LTR}\begin{tikzpicture}[baseline=-0.75ex,x=1em,y=1em]
 \fill[cFB] (0,0) circle (0.55);
 \node[white,inner sep=0,font=\usefont{TU}{lmss}{bx}{n}\small] at (0.04,-0.02) {f};
\end{tikzpicture}\end{LTR}}
\newcommand{\igicon}{\begin{LTR}\begin{tikzpicture}[baseline=-0.75ex,x=1em,y=1em]
 \shade[lower left=cIG1,lower right=cIG2,upper left=cIG2,upper right=cIG3,rounded corners=0.28em] (-0.55,-0.55) rectangle (0.55,0.55);
 \draw[white,line width=0.09em,rounded corners=0.18em] (-0.36,-0.36) rectangle (0.36,0.36);
 \draw[white,line width=0.09em] (0,0) circle (0.17);
 \fill[white] (0.22,0.22) circle (0.05);
\end{tikzpicture}\end{LTR}}
\newcommand{\tgicon}{\begin{LTR}\begin{tikzpicture}[baseline=-0.75ex,x=1em,y=1em]
 \fill[cTG] (0,0) circle (0.55);
 \fill[white] (-0.33,0.0) -- (0.32,0.27) -- (0.2,-0.3) -- (0.0,-0.15) -- (-0.08,-0.28) -- (-0.08,-0.08) -- cycle;
 \fill[cTG!70!black] (-0.08,-0.08) -- (0.25,0.2) -- (0.0,-0.15) -- (-0.08,-0.28) -- cycle;
\end{tikzpicture}\end{LTR}}
% page number: black box, white two-digit number, at the right of the social row
\newcommand{\PageBox}{\begin{LTR}\tikz[baseline=(p.base)]\node[fill=black,text=white,inner sep=2.5pt,minimum width=1.6em,minimum height=1.3em,font=\usefont{TU}{lmss}{bx}{n}\small](p){\ifnum\value{page}<10 0\fi\arabic{page}};\end{LTR}}
% visual left-to-right order: Telegram | Instagram | Facebook | page box
\newcommand{\SocialRow}{\begin{otherlanguage}{english}%
 \ar{\textbf{الأستاذ سوايسية هشام للرياضيات}}\ \tgicon\hfill
 \textbf{prof\_hicham\_math}\ \igicon\hfill
 \ar{\textbf{الأستاذ سوايسية هشام للرياضيات}}\ \fbicon\hspace{2.2em}\PageBox\end{otherlanguage}}

\pagestyle{fancy}
\fancyhf{}
\renewcommand{\headrulewidth}{0.6pt}
\renewcommand{\footrulewidth}{0.4pt}
% \SetHeader{right head: document title}{left head: level -- stream}
\newcommand{\SetHeader}[2]{%
 \fancyhead[C]{\parbox{\textwidth}{\ar{#1\hfill #2}}}%
 \fancyfoot[C]{\PageBox}}
\usepackage{tkz-tab}
\usepackage{multicol}
% variation / sign table (LTR)
\newenvironment{TAB}{\par\begin{center}\begin{LTR}\begin{tikzpicture}[scale=0.95,every node/.style={font=\small}]}{\end{tikzpicture}\end{LTR}\end{center}}
% grid + axes, clipped: \begin{GR}{xmin}{xmax}{ymin}{ymax}{unit cm}
\newenvironment{GR}[5]{\par\begin{center}\begin{LTR}\begin{tikzpicture}[x=#5cm,y=#5cm,font=\footnotesize]
 \draw[gray!25,very thin] (#1,#3) grid (#2,#4);
 \draw[-{Stealth[length=2mm]}] (#1,0)--(#2+0.3,0) node[below left]{$x$};
 \draw[-{Stealth[length=2mm]}] (0,#3)--(0,#4+0.3) node[below left]{$y$};
 \node[below left] at (0,0){$O$};
 \begin{scope}\clip (#1,#3) rectangle (#2,#4);}
{\end{scope}\end{tikzpicture}\end{LTR}\end{center}}
\newcommand{\Ticks}[2]{\foreach \t in {#1} \node[below,font=\scriptsize] at (\t,0){$\t$}; \foreach \t in {#2} \node[left,font=\scriptsize] at (0,\t){$\t$};}
\newcommand{\bx}[1]{\ensuremath{\colorbox{cDef!12}{$\displaystyle #1$}}}
\newcommand{\TwoCol}[4]{\par\begin{LTR}\noindent\begin{minipage}[t]{#1\textwidth}\selectlanguage{arabic}#2\end{minipage}\hfill\begin{minipage}[t]{#3\textwidth}\selectlanguage{arabic}#4\end{minipage}\end{LTR}\par}
\newcommand{\RL}[4]{\TwoCol{#3}{#4}{#1}{#2}}  % \RL{w1}{right content}{w2}{left content}
\newcommand{\ok}{$\textcolor{cRem}{\checkmark}$}
\newcommand{\limn}{\displaystyle\lim_{n\to+\infty}}
% sub-questions, each on its own line
\newenvironment{sq}{\begin{itemize}[leftmargin=2.4em,itemsep=3pt]}{\end{itemize}}
\newcommand{\qa}{\item[(أ)]}\newcommand{\qb}{\item[(ب)]}\newcommand{\qc}{\item[(جـ)]}\newcommand{\qd}{\item[(د)]}\newcommand{\qe}{\item[(هـ)]}
\newcommand{\pa}{\par\smallskip\textbf{(أ)} }\newcommand{\pb}{\par\smallskip\textbf{(ب)} }\newcommand{\pc}{\par\smallskip\textbf{(جـ)} }\newcommand{\pd}{\par\smallskip\textbf{(د)} }\newcommand{\pe}{\par\smallskip\textbf{(هـ)} }
\newcommand{\QQ}[1]{\par\smallskip\textbf{#1.} }
\newenvironment{tmr}{\stepcounter{exo}\begin{BOX}{cSec}{تمرين محلول \theexo}}{\end{BOX}}
\newcommand{\ABC}[3]{\\[2pt]\begin{tabular}{p{0.28\textwidth}p{0.28\textwidth}p{0.28\textwidth}}(جـ) \ #3 & (ب) \ #2 & (أ) \ #1\end{tabular}}

\newcommand{\nt}[1]{\hfill$\textcolor{cThm}{\mathbf{(#1)}}$}
\SetHeader{\small\bfseries\color{cSec}فرض المراقبة الذاتية رقم 01}{\small\color{cSec}السنة الثانية ثانوي -- تقني رياضي ورياضيات}
\begin{document}
\begin{center}
{\large\bfseries الجمهورية الجزائرية الديمقراطية الشعبية}\\[2pt]
{\bfseries وزارة التربية الوطنية}\\[2pt]
{\bfseries الديوان الوطني للتعليم والتكوين عن بعد}\\[2pt]
{\bfseries المركز الولائي سوق أهراس}
\end{center}
\noindent\textbf{السنة الدراسية: 2026/2027}\hfill\textbf{المستوى: سنة ثانية ثانوي شعبة تقني رياضي ورياضيات}
\par\smallskip
\begin{LTR}
\begin{tcolorbox}[enhanced,colback=cSec,colframe=cSec,arc=2mm,halign=center,boxrule=0pt,top=6pt,bottom=6pt,before upper={\selectlanguage{arabic}}]
{\color{white}\Large\bfseries فرض المراقبة الذاتية رقم 01}
\end{tcolorbox}
\end{LTR}

\begin{tamrin}{كثيرات الحدود (06 نقاط)}
نعتبر كثير الحدود $P$ المعرّف على $\R$ بـ: $P(x)=2x^3-3x^2-11x+6$.
\QQ{1} تحقّق أن $3$ جذر لكثير الحدود $P(x)$.
\QQ{2} عيّن الأعداد الحقيقية $a$ ، $b$ و $c$ بحيث من أجل كل عدد حقيقي $x$: $P(x)=(x-3)\left(ax^2+bx+c\right)$.
\QQ{3} حُلّ في $\R$ المعادلة: $P(x)=0$.
\QQ{4} ادرس إشارة $P(x)$ حسب قيم العدد الحقيقي $x$ ، ثم استنتج حلول المتراجحة: $P(x)\leqslant0$.
\QQ{5} استنتج في $\R$ حلول المعادلة: $2x^6-3x^4-11x^2+6=0$.
\end{tamrin}

\begin{tamrin}{الدوال: التفكيك واتجاه التغير والتناظر (06 نقاط)}
\textbf{I.} $f$ الدالة المعرّفة على $\R$ بـ: $f(x)=x^2-4x+3$ ، و $(C_f)$ تمثيلها البياني في معلم متعامد ومتجانس $\left(O;\vec i,\vec j\right)$.
\QQ{1} تحقّق أنه من أجل كل عدد حقيقي $x$: $f(x)=(x-2)^2-1$.
\QQ{2} عيّن دالتين مرجعيتين أو تآلفيتين $u$ و $v$ بحيث: $f=v\circ u$.
\QQ{3} استنتج اتجاه تغيّر الدالة $f$ على كل من المجالين $]-\infty;2]$ و $[2;+\infty[$ ، ثم شكّل جدول تغيّراتها.
\QQ{4} بيّن أن المستقيم ذا المعادلة $x=2$ محور تناظر للمنحنى $(C_f)$.
\par\medskip
\textbf{II.} $h$ الدالة المعرّفة على $\R-\{3\}$ بـ: $h(x)=\dfrac{2x-5}{x-3}$ ، و $(C_h)$ تمثيلها البياني.
\QQ{1} تحقّق أنه من أجل كل $x\neq3$: $h(x)=2+\dfrac1{x-3}$ ، ثم استنتج اتجاه تغيّر $h$ على كل من المجالين $]-\infty;3[$ و $]3;+\infty[$.
\QQ{2} بيّن أن النقطة $\Omega(3;2)$ مركز تناظر للمنحنى $(C_h)$.
\end{tamrin}

\begin{tamrin}{الاحتمالات (08 نقاط)}
يحتوي كيس على $5$ كريات متماثلة لا نفرّق بينها باللمس ، تحمل الأرقام: $0$ ، $1$ ، $1$ ، $2$ ، $2$.
نسحب عشوائياً كريتين من الكيس \textbf{على التوالي ودون إرجاع} ، ونسجّل الرقمين المحصل عليهما.
\QQ{1} مثّل مجموعة الإمكانيات بجدول ذي مدخلين ، ثم عيّن عددها.
\QQ{2} احسب احتمال كل حادثة من الحوادث التالية:
\begin{sq}
\qa $A$: «مجموع الرقمين يساوي $3$».
\qb $B$: «الكريتان تحملان نفس الرقم».
\qc $C$: «إحدى الكريتين على الأقل تحمل الرقم $0$».
\qd $A\cup C$.
\end{sq}
\QQ{3} ليكن $X$ المتغير العشوائي الذي يرفق بكل سحبة مجموع الرقمين المسجلين.
\begin{sq}
\qa عيّن القيم الممكنة للمتغير العشوائي $X$ ، ثم عرّف قانون احتماله.
\qb احسب الأمل الرياضياتي $E(X)$ ، والتباين $V(X)$ ، والانحراف المعياري $\sigma(X)$.
\end{sq}
\QQ{4} يدفع لاعب $m$ دينار للمشاركة في اللعبة ، ويتحصل على $10X$ دينار. ليكن $G$ ربحه الجبري.
\par عبّر عن $E(G)$ بدلالة $m$ ، ثم عيّن قيمة $m$ التي تجعل اللعبة عادلة.
\end{tamrin}
\begin{center}\textbf{\color{cSec}بالتوفيق}\end{center}

\newpage
\begin{LTR}
\begin{tcolorbox}[enhanced,colback=cSol,colframe=cSol,arc=2mm,halign=center,boxrule=0pt,top=6pt,bottom=6pt,before upper={\selectlanguage{arabic}}]
{\color{white}\Large\bfseries الحل النموذجي وسلم التنقيط}
\end{tcolorbox}
\end{LTR}
\setcounter{exo}{1}
\begin{hal}
\QQ{1} $P(3)=2(27)-3(9)-11(3)+6=54-27-33+6=\bx{0}$ ، إذن $3$ جذر لـ $P(x)$. \nt{0,5}
\QQ{2} بالنشر: $(x-3)\left(ax^2+bx+c\right)=ax^3+(b-3a)x^2+(c-3b)x-3c$.
\par بالمطابقة مع $P(x)$: $a=2$ ؛ $b-3a=-3$ ومنه $b=3$ ؛ $c-3b=-11$ ومنه $c=-2$ ؛ والتحقق: $-3c=6$.
\par إذن: \bx{P(x)=(x-3)\left(2x^2+3x-2\right)} \nt{1,5}
\QQ{3} $P(x)=0$ معناه $x-3=0$ أو $2x^2+3x-2=0$.
\par $\Delta=3^2-4(2)(-2)=25$ ، $x_1=\dfrac{-3-5}4=-2$ ، $x_2=\dfrac{-3+5}4=\dfrac12$.
\par مجموعة الحلول: \bx{S=\left\{-2;\tfrac12;3\right\}} \nt{1}
\QQ{4} $P(x)=(x-3)(x+2)(2x-1)$:
\begin{TAB}
\tkzTabInit[lgt=2.4,espcl=1.6]{$x$/0.8,$x+2$/0.7,$2x-1$/0.7,$x-3$/0.7,$P(x)$/0.8}{$-\infty$,$-2$,$\frac12$,$3$,$+\infty$}
\tkzTabLine{,-,z,+,t,+,t,+,}
\tkzTabLine{,-,t,-,z,+,t,+,}
\tkzTabLine{,-,t,-,t,-,z,+,}
\tkzTabLine{,-,z,+,z,-,z,+,}
\end{TAB}
حلول المتراجحة $P(x)\leqslant0$: \bx{S=]-\infty;-2]\cup\left[\tfrac12;3\right]} \nt{1,5}
\QQ{5} نضع $X=x^2$ ($X\geqslant0$) فتصبح المعادلة $2X^3-3X^2-11X+6=0$ أي $P(X)=0$.
\par ومنه $X=-2$ (مرفوض لأن $X\geqslant0$) أو $X=\dfrac12$ أو $X=3$.
\par $x^2=\dfrac12$ معناه $x=\pm\dfrac{\sqrt2}2$ ؛ $x^2=3$ معناه $x=\pm\sqrt3$.
\par مجموعة الحلول: \bx{S=\left\{-\sqrt3;-\tfrac{\sqrt2}2;\tfrac{\sqrt2}2;\sqrt3\right\}} \nt{1,5}
\end{hal}

\stepcounter{exo}
\begin{hal}
\textbf{I.}
\QQ{1} $(x-2)^2-1=x^2-4x+4-1=\bx{x^2-4x+3=f(x)}$ \nt{0,5}
\QQ{2} نضع $u(x)=x-2$ و $v(x)=x^2-1$ ، فيكون: $v\big(u(x)\big)=(x-2)^2-1=f(x)$ ، أي \bx{f=v\circ u} \nt{1}
\QQ{3} \textbf{على $]-\infty;2]$:} $u$ متزايدة و $u(x)\in]-\infty;0]$ ، و $v$ متناقصة على $]-\infty;0]$ ، إذن $f$ \textbf{متناقصة} على $]-\infty;2]$.
\par \textbf{على $[2;+\infty[$:} $u$ متزايدة و $u(x)\in[0;+\infty[$ ، و $v$ متزايدة على $[0;+\infty[$ ، إذن $f$ \textbf{متزايدة} على $[2;+\infty[$. \quad $f(2)=-1$.
\begin{TAB}
\tkzTabInit[lgt=2,espcl=2.6]{$x$/0.8,$f(x)$/1.6}{$-\infty$,$2$,$+\infty$}
\tkzTabVar{+/ ,-/$-1$,+/ }
\end{TAB}
\nt{1,5}
\QQ{4} من أجل كل $x\in\R$: $(4-x)\in\R$ و $f(4-x)=(4-x-2)^2-1=(2-x)^2-1=(x-2)^2-1=f(x)$.
\par إذن \bx{f(2\times2-x)=f(x)} ، والمستقيم $x=2$ محور تناظر لـ $(C_f)$. \nt{1}
\par\medskip
\textbf{II.}
\QQ{1} $2+\dfrac1{x-3}=\dfrac{2(x-3)+1}{x-3}=\dfrac{2x-5}{x-3}=h(x)$.
\par $h=w\circ g$ مع $g(x)=x-3$ (متزايدة) و $w(X)=2+\dfrac1X$ (متناقصة على كل من $]-\infty;0[$ و $]0;+\infty[$).
\par إذن $h$ \textbf{متناقصة} على كل من المجالين $]-\infty;3[$ و $]3;+\infty[$. \nt{1}
\QQ{2} من أجل كل $x\neq3$: $(6-x)\neq3$ و $h(6-x)+h(x)=2+\dfrac1{3-x}+2+\dfrac1{x-3}=4=2\times2$.
\par إذن \bx{h(2\times3-x)+h(x)=2\times2} والنقطة $\Omega(3;2)$ مركز تناظر لـ $(C_h)$. \nt{1}
\end{hal}

\stepcounter{exo}
\begin{hal}
\QQ{1} نرمز للكريات بـ: $0$ ، $1_a$ ، $1_b$ ، $2_a$ ، $2_b$. السحب على التوالي دون إرجاع: الترتيب مهم ولا تكرار ، فنشطب القطر. في كل خانة نكتب مجموع الرقمين (قيمة $X$):
\begin{center}\begin{LTR}{\renewcommand{\arraystretch}{1.5}\begin{tabular}{|c|c|c|c|c|c|}\hline
\rowcolor{cSec!12}$\text{I}\backslash\text{II}$ & $0$ & $1_a$ & $1_b$ & $2_a$ & $2_b$\\\hline
$0$ & \cellcolor{gray!30}$\times$ & \cellcolor{cRem!25}$1$ & \cellcolor{cRem!25}$1$ & \cellcolor{cRem!25}$2$ & \cellcolor{cRem!25}$2$\\\hline
$1_a$ & \cellcolor{cRem!25}$1$ & \cellcolor{gray!30}$\times$ & $2$ & \cellcolor{cEx!35}$3$ & \cellcolor{cEx!35}$3$\\\hline
$1_b$ & \cellcolor{cRem!25}$1$ & $2$ & \cellcolor{gray!30}$\times$ & \cellcolor{cEx!35}$3$ & \cellcolor{cEx!35}$3$\\\hline
$2_a$ & \cellcolor{cRem!25}$2$ & \cellcolor{cEx!35}$3$ & \cellcolor{cEx!35}$3$ & \cellcolor{gray!30}$\times$ & $4$\\\hline
$2_b$ & \cellcolor{cRem!25}$2$ & \cellcolor{cEx!35}$3$ & \cellcolor{cEx!35}$3$ & $4$ & \cellcolor{gray!30}$\times$\\\hline
\end{tabular}}\end{LTR}\end{center}
عدد الإمكانيات: \bx{5\times4=20} ، وهي متساوية الاحتمال. (البرتقالي: الحادثة $A$ ؛ الأخضر: الحادثة $C$.) \nt{1,5}
\QQ{2}
\pa $A$: الخانات البرتقالية: $8$ ، إذن \bx{P(A)=\tfrac{8}{20}=\tfrac25} \nt{0,5}
\pb $B$: الثنائيات $(1_a;1_b)$ ، $(1_b;1_a)$ ، $(2_a;2_b)$ ، $(2_b;2_a)$: \bx{P(B)=\tfrac4{20}=\tfrac15} \nt{0,5}
\pc $C$: الخانات الخضراء: $8$ ، إذن \bx{P(C)=\tfrac8{20}=\tfrac25} \nt{0,5}
\pd $A\cap C=\varnothing$ (مع الرقم $0$ لا يتجاوز المجموع $2$) ، إذن $P(A\cup C)=P(A)+P(C)-P(A\cap C)=\dfrac25+\dfrac25-0=\bx{\tfrac45}$ \nt{1}
\QQ{3}
\pa القيم الممكنة: $X\in\{1;2;3;4\}$. بعدّ الخانات: $P(X=1)=\frac4{20}$ ، $P(X=2)=\frac6{20}$ ، $P(X=3)=\frac8{20}$ ، $P(X=4)=\frac2{20}$:
\begin{center}\begin{LTR}{\renewcommand{\arraystretch}{1.6}\begin{tabular}{|C|c|c|c|c|C|}\hline
\rowcolor{cSec!12}$x_i$ & $1$ & $2$ & $3$ & $4$ & المجموع\\\hline
$P(X=x_i)$ & $\frac15$ & $\frac3{10}$ & $\frac25$ & $\frac1{10}$ & $1$\\\hline
\end{tabular}}\end{LTR}\end{center}
\nt{1,5}
\pb $E(X)=1\times\dfrac15+2\times\dfrac3{10}+3\times\dfrac25+4\times\dfrac1{10}=\dfrac{2+6+12+4}{10}=\bx{\tfrac{12}5=2{,}4}$ \nt{0,5}
\par $V(X)=\sum x_i^2P(X=x_i)-\big(E(X)\big)^2=\dfrac{2+12+36+16}{10}-\left(\dfrac{12}5\right)^2=\dfrac{33}5-\dfrac{144}{25}=\bx{\tfrac{21}{25}=0{,}84}$
\par $\sigma(X)=\sqrt{V(X)}=\bx{\tfrac{\sqrt{21}}5\approx0{,}92}$ \nt{1}
\QQ{4} $G=10X-m$ ، ومنه $E(G)=10E(X)-m=\bx{24-m}$.
\par اللعبة عادلة معناه $E(G)=0$ ، أي \bx{m=24} دينار. \nt{1}
\end{hal}

\begin{BOX}{cSec}{سلم التنقيط}
\begin{center}\begin{LTR}{\renewcommand{\arraystretch}{1.5}\begin{tabular}{|C|c|c|c|c|c|C|}\hline
\rowcolor{cSec!12}التمرين & \ar{س1} & \ar{س2} & \ar{س3} & \ar{س4} & \ar{س5} & المجموع\\\hline
الأول & $0{,}5$ & $1{,}5$ & $1$ & $1{,}5$ & $1{,}5$ & $06$\\\hline
الثاني & \multicolumn{5}{c|}{I: $0{,}5+1+1{,}5+1$ \quad II: $1+1$} & $06$\\\hline
الثالث & $1{,}5$ & $2{,}5$ & $3$ & $1$ & -- & $08$\\\hline
\rowcolor{cSec!12}المجموع & \multicolumn{5}{c|}{} & $20$\\\hline
\end{tabular}}\end{LTR}\end{center}
\end{BOX}
\end{document}
